\textbf{Proof.}
Fix \(G\) and a clique partition \(\mathcal P\), and write
\[
a_B=\frac{|B|}{n},\qquad B\in\mathcal P.
\tag{1}
\]

\textbf{A complete finite prior.}
For every sign array \(\xi=(\xi_{B,k}:B\in\mathcal P,k\in\{0,1\})\), define a complete potential-outcome schedule on \(H_{\mathcal P}\) as follows.  If \(i\in B\), put
\[
Y_i^{\xi}(z)=
\begin{cases}
\xi_{B,0},&z|_B\equiv0,\\
\xi_{B,1},&z_i=1\text{ and }z_j=0\text{ for all }j\in B\setminus\{i\},\\
0,&\text{otherwise}.
\end{cases}
\tag{2}
\]
By \(\text{def:clique-partition-subgraph}\), the closed neighborhood of \(i\) in \(H_{\mathcal P}\) is exactly \(B\).  Hence (2) depends only on \(z|_{\widetilde N_i(H_{\mathcal P})}\) and satisfies arbitrary neighborhood interference.  Moreover,
\[
Y_i^{\xi}(\mathbf 0)=\xi_{B,0},
\qquad
Y_i^{\xi}(\mathbf e_i)=\xi_{B,1}.
\tag{3}
\]
Consequently each atom satisfies both layer restrictions with equality:
\[
\frac1n\sum_{i\in V_n}Y_i^{\xi}(\mathbf0)^2
=\frac1n\sum_{B\in\mathcal P}|B|=1,
\qquad
\frac1n\sum_{i\in V_n}Y_i^{\xi}(\mathbf e_i)^2=1.
\tag{4}
\]
Thus every one of the \(4^{|\mathcal P|}\) schedules belongs to \(\mathcal M_2(H_{\mathcal P})\), by \(\text{def:general-published-q2-class}\).

Give these schedules the uniform prior, equivalently taking all signs mutually independent and Rademacher.  Equations (3) and \(\text{def:general-dte-primitives}\) give
\[
\tau_{\mathrm{DTE}}^{H_{\mathcal P}}(Y^{\xi})
=\sum_{B\in\mathcal P}a_B(\xi_{B,1}-\xi_{B,0}).
\tag{5}
\]

\textbf{Posterior obstruction for every assignment.}
Fix a binary assignment \(z\).  Within a block \(B\), there are three exhaustive cases.  If \(z|_B\equiv0\), the observations reveal \(\xi_{B,0}\) and not \(\xi_{B,1}\).  If exactly one unit is treated, its outcome reveals \(\xi_{B,1}\), while every other observed outcome in the block is the fixed value zero; thus \(\xi_{B,0}\) is not revealed.  If at least two units are treated, all observed outcomes in the block are the fixed value zero, so neither sign is revealed.  This includes singleton blocks: their two possible assignments reveal exactly one of their two signs.  Hence no assignment reveals both signs in any block.

All unrevealed signs remain conditionally independent Rademacher variables because the prior signs are independent and the non-target values in (2) are fixed independently of them.  Therefore
\[
\operatorname{Var}_{\xi}\!\left(
\tau_{\mathrm{DTE}}^{H_{\mathcal P}}(Y^{\xi})
\mid z,Y^{\xi}(z)
\right)
=\sum_{B\in\mathcal P}a_B^2N_B(z),
\tag{6}
\]
where \(N_B(z)\in\{1,2\}\) is the number of unrevealed signs in block \(B\).  In particular, (6) is at least \(\sum_Ba_B^2\).  By the clique structure and \(\mathrm{lem{:}published\text{-}conflict\text{-}graph\text{-}characterization}\), a maximal target-exposure observation chooses, independently block by block, either the all-control exposure or one isolated-treatment exposure.  It reveals exactly one sign in every block, so \(N_B(z)=1\) for every \(B\), and (6) is then exactly \(\sum_Ba_B^2\).

For squared loss and every estimator value \(t\), conditional completion of the square yields
\[
\mathbb E_{\xi}\!\left[(\tau-t)^2\mid z,Y^{\xi}(z)\right]
=\operatorname{Var}_{\xi}(\tau\mid z,Y^{\xi}(z))
+\left\{\mathbb E_{\xi}(\tau\mid z,Y^{\xi}(z))-t\right\}^2.
\tag{7}
\]
Thus every assignment law and every measurable estimator have Bayes risk at least \(\sum_Ba_B^2\).  Worst-case risk over \(\mathcal M_2(H_{\mathcal P})\) dominates this finite-prior Bayes risk.  Taking the infimum over the full design and estimator spaces in \(\text{def:general-dte-minimax-risk}\) gives
\[
R_2(H_{\mathcal P})\geq\sum_{B\in\mathcal P}a_B^2.
\tag{8}
\]

\textbf{Transfer and finite maximization.}
Because \(H_{\mathcal P}\subseteq G\) is spanning, \(\mathrm{lem{:}published\text{-}network\text{-}risk\text{-}monotonicity}\) gives
\[
R_2(G)\geq R_2(H_{\mathcal P}).
\tag{9}
\]
Combining (8) and (9) proves the certificate for every clique partition.  There are finitely many set partitions of the finite set \(V_n\), so maximizing over the subfamily whose blocks induce cliques proves the last display in the proposition.  This is a finite enumeration statement only and supplies no polynomial-time algorithm for the maximization.

As a smallest-instance check, if \(n=1\), the unique partition has one singleton block; every assignment reveals one of two independent signs and leaves the other unobserved, so (6) equals one exactly.  This agrees with the displayed certificate. \(\square\)